\section{Joint representation and orbit quotient}
\label{app:joint-representation}

The incidence automorphism group $G_A$ acts on the address Hilbert space
only through its base projection.  If
$\rho(\pi_B,\pi_T)=\pi_B$, then elements with $\pi_B=\mathrm{id}$ form the
kernel, and the faithfully represented group is
$G_B\simeq G_A/\ker\rho$.  Since the $S_k$ register action commutes with the
$G_B$ base action, complete reducibility gives
\begin{equation}
 \mathcal H_{\rm valid}\simeq
 \bigoplus_{\alpha,\beta}
 V_\alpha^{S_k}\otimes V_\beta^{G_B}\otimes
 \mathcal M_{\alpha\beta}.
 \label{eq:app-joint-decomposition}
\end{equation}
Every invariant Hamiltonian lies in the commutant and hence has the form in
Eq.~\eqref{eq:commutant-rewrite}.  The identity factors generate Schur
multiplicities; equality of eigenvalues belonging to different
$(\alpha,\beta)$ blocks is not enforced by this statement.

For completeness, let $\mathcal O$ and $\mathcal O'$ be computational-basis
orbits of the joint action and choose $x\in\mathcal O$.  Invariance gives
\begin{align}
 \langle\mathcal O'|H|\mathcal O\rangle
 &=\frac{1}{\sqrt{|\mathcal O||\mathcal O'|}}
 \sum_{x'\in\mathcal O}\sum_{y\in\mathcal O'}H_{yx'}\\
 &=\sqrt{\frac{|\mathcal O|}{|\mathcal O'|}}
 \sum_{y\in\mathcal O'}H_{yx},
\end{align}
which proves Eq.~\eqref{eq:orbit-quotient-rewrite}.  The orbit states are
therefore an exact orthonormal basis of the fixed space, not an effective
approximation.

After quotienting by register permutations, basis vectors are weak
multisets of size $k$ over the bases.  If $g\in G_B$ has $c_\ell(g)$ cycles
of length $\ell$, the trace of $g$ on
the symmetric tensor power $\operatorname{Sym}^k(\mathbb C^{n_B})$ is the
coefficient
\begin{equation}
 [z^k]\prod_{\ell\geq1}(1-z^\ell)^{-c_\ell(g)}.
\end{equation}
Averaging this character over $G_B$ projects onto the trivial irrep and
gives Eq.~\eqref{eq:burnside-rewrite}.  This counts
$\mathcal H_{\rm sym}$ only.  Whether $\ket u$ is cyclic for the restricted
endpoint algebra is a separate closure calculation.
